\documentclass[11pt,a4paper]{article}
\usepackage[utf8]{inputenc}
\usepackage{amsmath,amssymb,amsthm}
\usepackage{algorithm}
\usepackage{algpseudocode}
\usepackage{booktabs}
\usepackage{hyperref}
\usepackage{geometry}
\geometry{margin=1in}

\title{Teacher Forcing and Scheduled Sampling: Exposure Bias Mitigation, Curriculum Annealing, and Sequence Convergence}
\author{Team Scaibu \\ \small Email: \texttt{jaydeep.9664920749@gmail.com} \\ \small Architecture and Core Engine Team}
\date{\today}

\newtheorem{theorem}{Theorem}
\newtheorem{lemma}[theorem]{Lemma}
\newtheorem{proposition}[theorem]{Proposition}
\theoremstyle{definition}
\newtheorem{definition}{Definition}
\newtheorem{remark}{Remark}

\begin{document}

\maketitle

\begin{abstract}
Autoregressive sequence generation models suffer from exposure bias: during training under standard Teacher Forcing, the model is conditioned exclusively on true historical tokens $y_{1:t-1}^*$, whereas during inference, it must generate conditioned on its own potentially erroneous predictions $\hat{y}_{1:t-1}$. Scheduled Sampling bridges this train-test distribution mismatch by dynamically annealing a Bernoulli probability $\epsilon_i$ over training iterations, transitioning the decoder from full ground-truth guidance to self-generated token rollouts. This document formalizes the exposure bias phenomenon, proves distribution convergence under annealing schedules, details algorithmic pseudocode, and provides a verified step-by-step numerical example.
\end{abstract}

\section{Notation and Mathematical Preliminaries}

Let $\mathbf{y}^* = (y_1^*, \dots, y_T^*)$ denote the ground truth sequence, and let $\hat{y}_t \sim P_\theta(\cdot \mid \mathbf{x}, x_{1:t})$ denote the token predicted by the model at step $t$.

\begin{table}[h]
\centering
\caption{Mathematical Notation for Scheduled Sampling}
\begin{tabular}{lll}
\toprule
\textbf{Symbol} & \textbf{Domain} & \textbf{Semantic Description} \\
\midrule
$i$ & $\mathbb{N}_{\ge 0}$ & Global training iteration / epoch counter \\
$\epsilon_i$ & $[0, 1]$ & Probability of applying teacher forcing (feeding ground truth $y_{t-1}^*$) \\
$1 - \epsilon_i$ & $[0, 1]$ & Probability of feeding model prediction $\hat{y}_{t-1}$ as input $x_t$ \\
$k$ & $\mathbb{R}_{> 0}$ & Schedule decay curvature hyperparameter \\
$D_{\text{train}}, D_{\text{infer}}$ & Probability & Data distributions encountered during training and inference \\
$\epsilon_{\min}$ & $[0, 1]$ & Minimum floor for teacher forcing probability \\
\bottomrule
\end{tabular}
\end{table}

\section{Problem Definition and Core Concepts}

\subsection{3-Level Explanation}
\begin{enumerate}
    \item \textbf{Intuitive (High School level):} When learning to ride a bicycle, training wheels (teacher forcing) keep you from falling. But if you only ever ride with training wheels, the moment they are removed (inference), the first wobble makes you crash. Scheduled Sampling gradually lifts the training wheels higher and higher so you learn how to recover from wobbles.
    \item \textbf{Engineering Level:} Pure teacher forcing allows parallel GPU training in Transformers, but recurrent models suffer from cascading error drift when generating long sequences at test time. Scheduled sampling feeds a mixture of ground-truth tokens and argmax model predictions during training to train error-recovery paths.
    \item \textbf{Mathematical Level:} At training step $t$, the input token $x_t$ is sampled from a mixture distribution:
    \begin{equation}
    x_t \sim \epsilon_i \cdot \delta_{y_{t-1}^*} + (1 - \epsilon_i) \cdot P_\theta(\hat{y}_{t-1} \mid x_{<t}, \mathbf{x})
    \end{equation}
    under decay schedules such as the Inverse Sigmoid Decay $\epsilon_i = \frac{k}{k + \exp(i/k)}$.
\end{enumerate}

\subsection{5-Step Algorithmic Framework}
\begin{enumerate}
    \item \textbf{Decay Rate Computation:} Calculate current step probability $\epsilon_i = f(i; k)$.
    \item \textbf{Stochastic Bernoulli Trial:} At decoding step $t$, draw uniform random variable $r \sim \mathcal{U}(0, 1)$.
    \item \textbf{Input Token Assignment:} If $r < \epsilon_i$, set decoder input $x_t = y_{t-1}^*$; else set $x_t = \arg\max P(\hat{y}_{t-1})$.
    \item \textbf{Recurrent Forward Step:} Compute next hidden state $s_t = \text{RNN}(s_{t-1}, x_t)$ and predict target logits.
    \item \textbf{Loss Backpropagation:} Evaluate cross-entropy against true label $y_t^*$ and update weights.
\end{enumerate}

\section{Algorithm Specification}

\begin{algorithm}[H]
\caption{Scheduled Sampling Sequence Training Step}
\label{alg:scheduled_sampling}
\begin{algorithmic}[1]
\Require Ground truth $\mathbf{y}^* = (y_1^*, \dots, y_T^*)$, iteration $i$, parameter $k$, model $\theta$.
\Ensure Step cross-entropy loss $\mathcal{L}$.
\State $\epsilon_i \gets \frac{k}{k + \exp(i / k)}$
\State $s_0 \gets \text{InitState}()$
\State $x_1 \gets \langle \text{BOS} \rangle, \quad \mathcal{L} \gets 0$
\For{$t = 1$ \textbf{to} $T$}
    \State $s_t, \mathbf{z}_t \gets \text{DecoderStep}(s_{t-1}, x_t; \theta)$
    \State $\mathcal{L} \gets \mathcal{L} - \log \text{Softmax}(\mathbf{z}_t)_{y_t^*}$
    \State $\hat{y}_t \gets \arg\max_v \mathbf{z}_t[v]$
    \State $r \sim \mathcal{U}(0, 1)$
    \If{$r < \epsilon_i$}
        \State $x_{t+1} \gets y_t^*$ \Comment{Teacher forcing branch}
    \Else
        \State $x_{t+1} \gets \hat{y}_t$ \Comment{Self-generation branch}
    \EndIf
\EndFor
\State \Return $\mathcal{L}$
\end{algorithmic}
\end{algorithm}

\section{Theoretical Guarantees and Exposure Bias Bounds}

\begin{theorem}[Total Variation Bound on Exposure Discrepancy]
\label{thm:exposure_bias}
Let $P_{\text{infer}}(\mathbf{y})$ and $P_{\text{train}}(\mathbf{y}; \epsilon)$ denote the generated prefix distributions. The Total Variation distance between the inference and training input distributions satisfies:
\begin{equation}
\text{TV}(D_{\text{train}}(\epsilon_i), D_{\text{infer}}) \le \epsilon_i \cdot \sum_{t=1}^T \mathbb{E}_{x_{<t}}[1 - P_\theta(y_t^* \mid x_{<t})]
\end{equation}
As $\epsilon_i \to 0$, the discrepancy vanishes: $\lim_{i \to \infty} \text{TV}(D_{\text{train}}, D_{\text{infer}}) = 0$.
\end{theorem}

\begin{proof}
At step $t$, the probability of divergence between training and inference inputs is upper-bounded by $\epsilon_i \cdot \mathbb{P}(\hat{y}_{t-1} \ne y_{t-1}^*)$. Summing across all $T$ sequence steps via the union bound gives the stated total variation bound.
\end{proof}

\section{Complexity Analysis}

\begin{itemize}
    \item \textbf{Time Complexity:} $\mathcal{O}(T \cdot d_h^2)$ operations per sequence (requires sequential decoding during training when $\epsilon_i < 1$).
    \item \textbf{Space Complexity:} $\mathcal{O}(T \cdot d_h)$ memory for gradient backpropagation.
\end{itemize}

\section{Exactness and Error Analysis}

While Scheduled Sampling is statistically non-rigorous (the gradients through the discrete $\arg\max$ token choice are ignored), in practice it provides substantial empirical BLEU/ROUGE improvements over pure teacher forcing.

\section{Worked Numerical Example}

Let $k = 100$, iteration $i = 100$.
\begin{itemize}
    \item Scaled ratio: $i / k = 100 / 100 = 1.0$.
    \item $\exp(1.0) \approx 2.71828$.
    \item $\epsilon_{100} = \frac{100}{100 + 2.71828} = \frac{100}{102.71828} \approx 0.9735$.
\end{itemize}
At iteration $i=500$:
\begin{itemize}
    \item $i / k = 5.0 \implies \exp(5.0) \approx 148.413$.
    \item $\epsilon_{500} = \frac{100}{100 + 148.413} = \frac{100}{248.413} \approx 0.4025$.
\end{itemize}

\section{Numerical Considerations and Edge Cases}

\begin{itemize}
    \item \textbf{Transformers vs RNNs:} In Transformers, teacher forcing enables $\mathcal{O}(1)$ parallel forward steps via causal masks. Scheduled sampling breaks this parallelism, requiring $\mathcal{O}(T)$ sequential passes.
\end{itemize}

\section{References}
\begin{enumerate}
    \item Bengio, S., Vinyals, O., Jaitly, N., & Shazeer, N. (2015). Scheduled sampling for sequence prediction with recurrent neural networks. \emph{NeurIPS}, 1171--1179.
    \item Lamb, A. M., Goyal, A. G. A. P., Zhang, Y., Zhang, S., Courville, A. C., & Bengio, Y. (2016). Professor forcing: Designing networks for continuous test-time inputs. \emph{NeurIPS}, 2091--2099.
\end{enumerate}

\end{document}
